% ==================== 中英文字体、数学字体和默认行距 ====================

%%%%%%%% 字体、段落 %%%%%%%%

% 采用 Windows 字体文件以兼容 macOS 和 Linux
\setCJKfamilyfont{simsun}{simsun.ttc}
\setCJKfamilyfont{simhei}{simhei.ttf}
\setCJKfamilyfont{simkai}{simkai.ttf}

\renewcommand{\songti}{\CJKfamily{simsun}} % 宋体
\renewcommand{\heiti}{\CJKfamily{simhei}}  % 黑体
\renewcommand{\kaishu}{\CJKfamily{simkai}} % 楷体

% 英文默认字体
\setmainfont{Times New Roman}
\setsansfont{Arial}
% 中文默认字体
\setCJKmainfont[AutoFakeBold=2.17,ItalicFont=simkai.ttf]{simsun.ttc}
\setCJKsansfont{simhei.ttf}
% 这样设置之后，针对不同的中英字体效果如下：
% \rm：中文宋体，英文 Times New Roman
% \sf：中文黑体，英文 Arial
% \bf：中文宋体加粗，英文 Times New Roman 加粗
% \it：中文楷体，英文 Times New Roman 斜体
% \bf\songti：中文宋体加粗，英文 Times New Roman 加粗

% 通用格式显式配对中西文字体；声明页仍可单独指定各自的专项字体。
\newcommand{\xmu@songti}{\normalfont\rmfamily\songti}
\newcommand{\xmu@heiti}{\normalfont\sffamily\heiti}

% 数学加粗斜体
\DeclareMathAlphabet{\mathbfit}{T1}{ptm}{b}{it}
\renewcommand{\boldsymbol}[1]{\mathbfit{#1}}
% URL 字体
\urlstyle{same}
% 默认行距
\renewcommand{\baselinestretch}{1.5}
